\section*{Bab \ref{ch:g7:central} --- Simetri Pusat dan Jajargenjang}

\begin{solution}{exo:g7:central:1}
Bayangan setengah putaran: $M'$ berada empat petak ke \emph{kiri} dari
$O$ dan satu petak ke \emph{bawah} (kedua arahnya terbalik). Bayangan
pencerminan terhadap garis tegaknya: empat petak ke kiri tetapi tetap
satu petak ke \emph{atas}. Kedua bayangannya berbeda: keduanya
bayangan cermin tegak satu sama lain.
\end{solution}

\begin{solution}{exo:g7:central:2}
$A'B' = AB$ (setengah putaran mempertahankan panjang) dan
$(A'B') \parallel (AB)$ (bayangan sebuah garis adalah garis yang
\omterm{def:g6:lines:perp}{sejajar}, \cref{prop:g7:central:preserved}).
\end{solution}

\begin{solution}{exo:g7:central:3}
Angka kalkulator yang punya \omterm{def:g7:central:center}{pusat simetri}: $0$, $2$, $5$, $8$ (dan
$1$ kalau digambar sebagai batang polos). Hurufnya: H, N, S, Z, O
punya; A dan T tidak (keduanya punya sumbu, bukan pusat).
\end{solution}

\begin{solution}{exo:g7:central:4}
Sisi yang berhadapan sama panjang: $MN = KL = 7$ cm dan
$NK = LM = 4$ cm. Sudut yang berhadapan sama besar:
$\widehat M = \widehat K = 115^\circ$. Diagonalnya saling memotong di
titik tengah bersamanya, dan titik itu adalah $O$: jadi titik tengah
$[LN]$ adalah $O$.
\end{solution}

\begin{solution}{exo:g7:central:5}
Gambarlah \omterm{def:g6:lines:objects}{ruas garis} $[AC]$ sepanjang $8$ cm, tandai titik tengahnya
$O$, gambarlah lewat $O$ sebuah garis yang membentuk sudut $60^\circ$
dengan $(AC)$, lalu letakkan $B$ dan $D$ padanya berjarak $2.5$ cm
dari $O$ di kedua sisinya. Hubungkan $A$, $B$, $C$, $D$: diagonalnya
berpotongan di titik tengahnya, jadi $ABCD$ adalah \omterm{def:g7:central:parallelogram}{jajargenjang}
menurut definisinya.
\end{solution}

\begin{solution}{exo:g7:central:6}
Sudut yang berurutan berpelurus: di sebelah $115^\circ$ duduk
$180 - 115 = 65^\circ$. Sudut yang berhadapan sama besar, jadi keempat
sudutnya adalah $115^\circ$, $65^\circ$, $115^\circ$, $65^\circ$
(\omterm{def:g1:addition:def}{jumlahnya} $360^\circ$).
\end{solution}

\begin{solution}{exo:g7:central:7}
Dari $A(1,1)$ ke $O(2.5, 2.5)$: $1.5$ ke kanan dan $1.5$ ke atas;
dengan melanjutkan langkah yang sama: $A'(4, 4)$. Dari $B(4,2)$ ke
$O$: $1.5$ ke kiri dan $0.5$ ke atas; dengan melanjutkan: $B'(1, 3)$.
\end{solution}

\begin{solution}{exo:g7:central:8}
Tidak. Trapesium sama kaki punya dua sisi miring yang sama panjang
tanpa menjadi \omterm{def:g7:central:parallelogram}{jajargenjang} (kedua sisi yang \omterm{def:g6:lines:perp}{sejajar} berbeda panjang).
Yang mencukupi adalah kriteria 3 pada
\cref{thm:g7:central:recognize}: dua sisi yang berhadapan sama panjang
\emph{dan \omterm{def:g6:lines:perp}{sejajar}}.
\end{solution}

\begin{solution}{exo:g7:central:9}
$O$ adalah titik tengah $[BC]$ menurut pilihan kita, dan ia titik
tengah $[AA']$ menurut pelukisan titik simetrisnya. \omterm{def:g4:geometry:polygon}{Segi empat}
$ABA'C$ punya diagonal $[AA']$ dan $[BC]$ yang berbagi titik tengah
$O$: kriteria 1 (definisinya) menjadikannya \omterm{def:g7:central:parallelogram}{jajargenjang} tanpa biaya
tambahan.
\end{solution}

\begin{solution}{exo:g7:central:10}
``Sudut berhadapan sama besar $\Rightarrow$ \omterm{def:g7:central:parallelogram}{jajargenjang}'':
\emph{benar} (ini kriteria pengenalan yang lain; dengan \omterm{def:g1:addition:def}{jumlah} sudut
$360^\circ$, pasangan berhadapan yang sama besar memaksa sudut
berurutan menjadi berpelurus, jadi sisinya \omterm{def:g6:lines:perp}{sejajar} menurut
\cref{thm:g7:triangles:equalities}).

``Diagonal yang \omterm{def:g6:lines:perp}{tegak lurus} $\Rightarrow$ \omterm{def:g6:shapes:quadrilaterals}{belah ketupat}'':
\emph{benar} untuk \omterm{def:g7:central:parallelogram}{jajargenjang} (diagonalnya memang sudah berpotongan
di titik tengahnya).

``Satu \omterm{def:g3:shapes:rightangle}{sudut siku-siku} $\Rightarrow$ persegi panjang'': \emph{benar}
--- sudut yang berurutan berpelurus, jadi keempatnya menjadi siku-siku.
\end{solution}

\begin{solution}{exo:g7:central:11}
Apa pun \omterm{def:g4:geometry:polygon}{segi empatnya} --- bahkan yang sangat tak beraturan --- titik
tengah $I$, $J$, $K$, $L$ selalu membentuk \emph{\omterm{def:g7:central:parallelogram}{jajargenjang}}.
Buktinya dengan teorema titik tengah ada pada \cref{ch:g8:midpoints}
(\cref{exo:g8:midpoints:9}).
\end{solution}

\begin{solution}{pb:g7:central:1}
\textbf{1.} $M(3, 1) \to (-3, -1)$; $N(-2, 4) \to (2, -4)$;
$P(0, -5) \to (0, 5)$. Aturannya: setengah putaran terhadap titik asal
mengirim $(x, y)$ ke $(-x, -y)$ --- kedua koordinatnya berganti tanda.

\textbf{2.} Dari $M(5, 3)$: tiga petak ke kanan dari $C$ menjadi
tiga petak ke kiri, dua ke atas menjadi dua ke bawah: bayangannya
$(-1, -1)$. Periksa resepnya: dua kali pusatnya dikurangi titiknya,
$(2 \times 2 - 5,\ 2 \times 1 - 3) = (-1, -1)$.

\textbf{3.} Setengah putaran terhadap pusat kartunya menukar
lambangnya berpasangan. Dengan lambang berjumlah \omterm{def:g3:numbers:evenodd}{ganjil}, satu lambang
tidak punya pasangan: ia harus menjadi bayangannya sendiri, dan
satu-satunya titik yang menjadi bayangannya sendiri adalah pusatnya.
Karena itu lambang tengah pada kartu 3, kartu 5, kartu
7\,\dots\ duduk tepat di pusat (lihatlah kartu sungguhan: memang benar).

\textbf{4.} Terhadap $O_1(2,0)$: $A(1,1) \to (3,-1)$; terhadap
$O_2(5,0)$: $(3,-1) \to (7,1)$. Awalnya $A(1,1)$, akhirnya
$(7,1)$: bergeser $6$ ke kanan, tidak terbalik. Sama untuk
$B(0,3) \to (4,-3) \to (6,3)$: bergeser $6$ ke kanan. Kedua setengah
putaran itu setara dengan sebuah \emph{translasi}.

\textbf{5.} Geserannya $6$, dan $O_1 O_2 = 3$: translasinya menempuh
\emph{dua kali} jarak antara kedua pusatnya (dengan arah dari $O_1$ ke
$O_2$) --- persis seperti dua cermin \omterm{def:g6:lines:perp}{sejajar} yang menggeser sejauh dua
kali celahnya pada \cref{pb:g6:symmetry:1}.

\textbf{6.} Setengah putaran terhadap $O$ mengirim \omterm{def:g7:central:parallelogram}{jajargenjangnya} ke
dirinya sendiri (itulah arti ``\omterm{def:g7:central:center}{pusat simetri}'' baginya,
\cref{thm:g7:central:properties}). Garis melalui $O$ juga terkirim ke
dirinya sendiri, jadi kedua titik tempat garis itu memotong tepinya
saling bertukar: keduanya \omterm{def:g7:central:def}{simetris} terhadap $O$. Kedua potongan
\omterm{def:g7:central:parallelogram}{jajargenjangnya} di sisi mana pun dari garis itu juga bertukar, jadi
keduanya salinan yang persis (\cref{prop:g7:central:preserved}) ---
luasnya sama.

\textbf{7.} Potonglah sepanjang garis yang melalui \emph{kedua
pusatnya}: pusat persegi panjang kuenya dan pusat lubangnya. Menurut
pertanyaan 6 yang dipakai pada persegi panjang kuenya, potongan itu
membelah dua kuenya; dipakai pada persegi panjang lubangnya (garisnya
melalui pusatnya juga), ia membelah dua lubangnya. Tiap bagian
memperoleh setengah kue dikurangi setengah lubang: adil sempurna.

\textbf{8.} $D = (0 + 8 - 6,\ 0 + 5 - 1) = (2, 4)$ (diagonalnya
berbagi titik tengah, jadi $D$ adalah bayangan setengah putaran dari
$B$ terhadap pusatnya). Pusatnya: titik tengah $[AC]$, yaitu
$O(4,\ 2.5)$.

\textbf{9.} Pada \omterm{def:g7:central:parallelogram}{jajargenjang} $ABA'C$, sisi yang berhadapan sama
panjang: $BA' = AC$. Ketaksamaan segitiga pada $ABA'$ memberi
$AA' < AB + BA' = AB + AC$. Karena $O$ adalah titik tengah
$[AA']$, berlaku $AA' = 2\,AO$, jadi $2\,AO < AB + AC$, yaitu
$AO < \frac{AB + AC}{2}$.

\textbf{10.} Batas atasnya: $AO < \frac{5 + 7}{2} = 6$ cm. Batas
bawahnya: pada $ABA'$, $AA' > BA' - AB = 7 - 5 = 2$, jadi
$AO > 1$ cm. Garis berat dari $A$ terletak benar-benar di antara $1$
dan $6$ cm.

\textbf{11.} \omterm{def:g6:lines:objects}{Ruas garis}: ya, titik tengahnya. Garis: ya --- setiap
titiknya adalah pusat. \omterm{def:g6:shapes:triangles}{Segitiga sama sisi}: tidak ada. Persegi:
ya, titik potong diagonalnya. Lingkaran: ya, pusatnya.
S dan Z: ya, titik tengahnya (balikkan halamannya: keduanya terbaca
sama).

\textbf{12.} Setengah putaran itu akan memasangkan ketiga \omterm{def:g5:solids:def}{titik
sudutnya} satu sama lain; karena tiga itu \omterm{def:g3:numbers:evenodd}{ganjil}, satu \omterm{def:g5:solids:def}{titik sudut} $A$
akan menjadi bayangannya sendiri, sehingga $A$ terpaksa menjadi pusat
setengah putarannya. Kedua \omterm{def:g5:solids:def}{titik sudut} lain $B$ dan $C$ lalu akan
\omterm{def:g7:central:def}{simetris} terhadap $A$ --- jadi $A$ akan menjadi titik tengah $[BC]$,
yang menaruh ketiga \omterm{def:g5:solids:def}{titik sudutnya} pada satu garis. Segitiga tidak
punya tiga \omterm{def:g5:solids:def}{titik sudut} yang segaris: kontradiksi. Pusatnya tidak ada.

\textbf{13.} Alasan keganjilan yang sama: setengah putaran yang
mempertahankan \omterm{def:g4:geometry:polygon}{segi banyaknya} memasangkan \omterm{def:g5:solids:def}{titik sudutnya}; \omterm{def:g5:solids:def}{titik sudut}
yang berjumlah \omterm{def:g3:numbers:evenodd}{ganjil} memaksa adanya \omterm{def:g5:solids:def}{titik sudut} tetap, yang harus
menjadi pusatnya sekaligus titik tengah \omterm{def:g6:lines:objects}{ruas garis} yang menghubungkan
dua \omterm{def:g5:solids:def}{titik sudut} lain --- tiga \omterm{def:g5:solids:def}{titik sudut} yang segaris, mustahil pada
\omterm{def:g4:geometry:polygon}{segi banyak}\,\dots\ jadi tidak ada \omterm{def:g4:geometry:polygon}{segi banyak} dengan \omterm{def:g5:solids:def}{titik sudut}
berjumlah \omterm{def:g3:numbers:evenodd}{ganjil} yang punya \omterm{def:g7:central:center}{pusat simetri}.

\textbf{14.} Melakukan setengah putaran terhadap $O_1$, lalu terhadap
$O_2$, mengirim bangunnya ke dirinya sendiri dua kali --- jadi
gabungan keduanya juga begitu. Menurut pertanyaan 5, gabungan itu
adalah translasi sejauh dua kali jarak $O_1 O_2$, yang bukan \omterm{def:g1:counting:zero}{nol}
karena kedua pusatnya berbeda. Bangun yang muat pada selembar kertas
tidak dapat sama dengan dirinya yang tergeser sejauh jarak tetap yang
bukan \omterm{def:g1:counting:zero}{nol} (geserlah cukup sering dan ia meninggalkan kertasnya sama
sekali). Jadi bangun yang terbatas punya paling banyak satu \omterm{def:g7:central:center}{pusat
simetri}.

\textbf{15.} Pada jalur tanpa ujung itu, setengah putaran terhadap
titik yang berada di tengah antara satu jejak kiri dan jejak kanan
berikutnya mengirim polanya ke dirinya sendiri; titik yang berjarak
satu langkah penuh lebih jauh juga begitu: dua pusat yang berbeda.
Menggabungkan kedua setengah putaran itu memberi translasi sejauh dua
kali setengah langkahnya --- persis geseran satu langkah yang tampak
memetakan pita hias tak berhingga itu ke dirinya sendiri, seperti yang
diramalkan pertanyaan 14. Pola yang tak terbatas hidup dengan aturan
yang berbeda: itulah matematika kertas dinding.

\end{solution}
